% BHU PYQ -- Manually transcribed & error-corrected
% Paper: MTB-503 : Programming in C
% Exam: B.A./B.Sc. (Hons) Semester V Examination, 2022-23

\documentclass[11pt,a4paper]{article}
\usepackage[a4paper,top=2.5cm,bottom=2.5cm,left=2.8cm,right=2.8cm,headheight=14pt]{geometry}
\usepackage{amsmath,amssymb}
\usepackage[T1]{fontenc}\usepackage[utf8]{inputenc}
\usepackage{lmodern,microtype,fancyhdr,enumitem,hyperref,lastpage,parskip}
\usepackage{listings}
\lstset{basicstyle=\small tfamily, breaklines=true, frame=single, xleftmargin=1cm}

\hypersetup{colorlinks=true,urlcolor=black,linkcolor=black,
  pdftitle={MTB-503: Programming in C -- BHU B.A./B.Sc. Sem V 2022-23}}
\pagestyle{fancy}\fancyhf{}
\renewcommand{\headrulewidth}{0.4pt}\renewcommand{\footrulewidth}{0.4pt}
\fancyhead[L]{\small   \textit{Banaras Hindu University}}
\fancyhead[R]{\small   \textit{MTB-503: Programming in C}}
\fancyfoot[C]{\small Page  \thepage\ of \pageref{LastPage}}


\newlist{parts}{enumerate}{2}
\setlist[parts,1]{label=(\alph*),leftmargin=*,itemsep=5pt,topsep=3pt}
\setlist[parts,2]{label=(\roman*),leftmargin=*,itemsep=3pt,topsep=2pt}

\newcommand{\pts}[1]{\hfill   \textit{[#1]}}


\begin{document}

\begin{center}
  {\large\bfseries BANARAS HINDU UNIVERSITY}\\[2pt]
  {
\normalsize B.A./B.Sc. (Hons) --- Semester V Examination, 2022-23}\\[6pt]
  \rule{0.8\linewidth}{0.5pt}\\[6pt]
  {\large\bfseries Paper No. MTB-503: Programming in `C'}\\[4pt]
  {
\normalsize Time: 3 Hours \hfill Full Marks: 50}
\end{center}
\rule{\linewidth}{0.4pt}
\smallskip



\noindent   \textit{Instructions: Answer   \textbf{five} questions including Question~1 which is compulsory. Figures in right-hand margin indicate marks.}
\bigskip




\noindent   \textbf{Question 1.}

\begin{parts}
  \item What is the use of   \texttt{printf()} and   \texttt{scanf()} functions? \pts{1}
  \item Name four header files in the `C' programming language. \pts{1}
  \item Determine the hierarchy of operations and evaluate the following expression, assuming   \texttt{k} is a   \texttt{float} variable:   \texttt{k = 3/2*4+3/8}. \pts{1}
  \item What is the use of the function   \texttt{toupper()}? Explain with an example code. \pts{1}
  \item If   \texttt{a = 10},   \texttt{b = 12} and   \texttt{c = 0}, find the values of: (i)   \texttt{a != 6 \&\& b > 5}; (ii)   \texttt{!(a > 5 \&\& c)}. \pts{1}
  \item What are the basic data types in `C'? \pts{1}
  \item What is the output of the following program?
  
\begin{lstlisting}
main() {
    int x = 101, y = 105;
    printf("%d", (x > y) ? x : y);
}
  \end{lstlisting} \pts{1}
  \item What is the difference between   \texttt{++a} and   \texttt{a++}? \pts{1}
  \item Describe the difference between   \texttt{=} and   \texttt{==} operators in `C'. \pts{1}
  \item Write the equivalent   \texttt{while} loop code for:
  
\begin{lstlisting}
for (int i = 0; i < 5; i++)
    printf("%d", i);
  \end{lstlisting} \pts{1}
\end{parts}

\medskip\rule{\linewidth}{0.2pt}

\medskip


\noindent   \textbf{Question 2.}

\begin{parts}
  \item Write a brief history of the `C' programming language. \pts{5}
  \item Write a `C' program to find the Greatest Common Divisor (GCD) of two natural numbers. \pts{5}
\end{parts}

\medskip\rule{\linewidth}{0.2pt}

\medskip


\noindent   \textbf{Question 3.}

\begin{parts}
  \item What are the differences between local and global variables in `C' with reference to declaration, visibility, and lifetime? \pts{5}
  \item Write a general `C' program to find the sum of the series $x + \dfrac{x^3}{3!} + \dfrac{x^5}{5!} + \cdots + \dfrac{x^{2n-1}}{(2n-1)!}$ and print it on the screen. \pts{5}
\end{parts}

\medskip\rule{\linewidth}{0.2pt}

\medskip


\noindent   \textbf{Question 4.}

\begin{parts}
  \item What are entry-controlled and exit-controlled loops? Give examples. Using a   \texttt{while} loop, write a `C' program to evaluate $y = x^n$ where $n$ is a non-negative integer. \pts{5}
  \item Write a `C' program to find the length of a string and count the number of vowels in it. \pts{5}
\end{parts}

\medskip\rule{\linewidth}{0.2pt}

\medskip


\noindent   \textbf{Question 5.}

\begin{parts}
  \item Write a `C' program to print the multiplication table from $1 \times 1$ to $12  \times 10$. \pts{5}
  \item What is an array? Write the declaration, initialization and memory allocation of a one-dimensional array. \pts{5}
\end{parts}

\medskip\rule{\linewidth}{0.2pt}

\medskip


\noindent   \textbf{Question 6.}

\begin{parts}
  \item Write a `C' program to compute and print the sum of all elements stored in an array using a pointer. \pts{5}
  \item Write a `C' program to: (i) print the word   \texttt{MATHEMATICS}; (ii) count the occurrences and location of the letter `A' in that word. \pts{5}
\end{parts}

\medskip\rule{\linewidth}{0.2pt}

\medskip


\noindent   \textbf{Question 7.}

\begin{parts}
  \item Write a `C' program that uses functions and structures to perform: (i) reading a complex number; (ii) writing a complex number; (iii) addition of two complex numbers; (iv) multiplication of two complex numbers. \pts{5}
  \item Write a `C' program to read a series of three-digit integers from a file named   \texttt{DATA}, then write all odd numbers to a file called   \texttt{ODD} and all even numbers to a file called   \texttt{EVEN}. \pts{5}
\end{parts}

\vfill\rule{\linewidth}{0.4pt}

\begin{center}\small
  \href{https://bhu-exam.vercel.app/}{Click here to visit our website}\\[4pt]
  \href{https://me-aryan.vercel.app/}{Click here to visit our contributor}\\[4pt]
  \href{https://quashan-me.vercel.app/}{Click here to visit our contributor}
\end{center}
\end{document}
